problem:  
**Conjecture (Sublinear Liouville–Cone Uniqueness Conjecture).**  
Let \( (M^n,g) \) be a complete noncompact Riemannian manifold satisfying  
\[
\mathrm{Ric}(x) \ge -\frac{C}{(1+r(x))^2}, \qquad 
\operatorname{Vol}(B(x_0,r)) \ge c\, r^n,
\]
for some constants \(C,c>0\), where \(r(x) = d(x,x_0)\).  
Assume that \(M\) satisfies the *sublinear Liouville property*: every harmonic function \(u\colon M\to\mathbb{R}\) with 
\[
|u(x)| = o(r(x)) \quad \text{as } r(x)\to\infty
\]
is constant.  

**Question (precise form):**  
Does it follow that \(M\) has a *unique tangent cone at infinity* in the measured Gromov–Hausdorff sense?  

Equivalently, is it true that the disappearance of all nontrivial sublinear harmonic functions enforces global geometric uniqueness of asymptotic cones under the borderline curvature condition \(\mathrm{Ric}\ge -C(1+r)^{-2}\) and noncollapsing volume growth?  

Why is it a "good" problem:  
This conjecture presents a sharp and conceptually minimal bridge between *analytic rigidity* (the extreme Liouville property) and *geometric rigidity* (uniqueness of tangent cones) under the critical Ricci curvature decay where standard nonnegative Ricci techniques no longer directly apply.  

In the \(\mathrm{Ric}\ge0\) case, Cheeger–Colding and Colding–Minicozzi established a deep relationship between harmonic function spaces and the geometry at infinity: variation in tangent cones corresponds to abundance of polynomial-growth harmonic functions.  The present conjecture asks whether the *complete suppression* of sublinear harmonic functions suffices to rigidify the asymptotic geometry when curvature only barely fails to be nonnegative.  

Its simplicity—“Does sublinear Liouville imply unique cone under borderline curvature decay?”—conceals profound depth.  Resolving it would require new analytic tools controlling harmonic functions at infinity under near-Ricci-flat conditions, potentially producing a unifying mechanism linking Liouville properties, volume noncollapsing, and metric–measure asymptotics.  A positive answer would extend Colding–Minicozzi–Cheeger–Colding rigidity to a natural boundary of curvature control; a counterexample would reveal the precise limits of analytic rigidity as a detector of global geometry.